\chapter{Introduction}

Financial markets are fundamentally information-processing systems. 
Asset prices reflect expectations about future cash flows, risk, and growth opportunities, all of which depend on the continuous arrival of new information. 
Among the most economically significant and recurring information events in equity markets are corporate earnings announcements. 
Quarterly and annual financial reports disclose updated accounting performance, managerial assessments, and forward-looking guidance, often triggering substantial stock price adjustments. 
Understanding the short-term effect of financial reports on stock price movements remains a central problem in financial economics and, increasingly, in machine learning-driven asset pricing research.

The informational role of earnings announcements has been documented for more than five decades. 
In their seminal study, Ray Ball and Philip Brown\cite{an-empirical-evaluation-of-accounting-income-numbers} provide early empirical evidence that accounting earnings convey value-relevant information that is rapidly incorporated into stock prices. 
Their findings established the foundation for capital markets research in accounting and demonstrated that financial disclosures are not merely regulatory formalities but economically meaningful signals.

Subsequent literature has consistently shown that earnings announcements are associated with abnormal returns, elevated volatility, and abnormal trading volume. 
For example, Beaver\cite{information-content-of-annual-earnings-announcements} documents significant increases in trading activity around earnings releases, indicating that new information is actively processed by market participants. 
Later research identifies systematic return patterns following earnings surprises. 
Notably, Victor Bernard and Jacob Thomas\cite{pead} document post-earnings announcement drift (PEAD), showing that stock prices continue to drift in the direction of earnings surprises for weeks after the announcement. 
This phenomenon suggests that price adjustment to earnings information may be gradual rather than instantaneous.

More recent research continues to confirm the persistence of announcement-related anomalies, highlighting the ongoing relevance of earnings-based predictability (e.g., Livnat \& Mendenhall\cite{compare-pead-analysts-vs-time-series}). 
The magnitude and consistency of these findings explain why institutional investors and quantitative funds devote substantial resources to analyzing earnings releases in real time.

The theoretical benchmark for evaluating price reactions to financial reports is the Efficient Market Hypothesis (EMH), formalized by Eugene Fama\cite{efficient-capital-markets}. 
Under the semi-strong form of the EMH, publicly available information, such as financial statements, should be immediately and fully reflected in stock prices. 
If markets are semi-strong efficient, systematic short-term predictability following earnings announcements should not persist.

However, empirical evidence of underreaction, post-announcement drift, and earnings momentum challenges the strictest interpretation of semi-strong efficiency. 
Fama\cite{market-efficiency-long-term-returns} acknowledges that anomalies exist but argues that many may be artifacts of data-snooping or model misspecification. 
Nonetheless, the persistence of earnings-related predictability across time periods and markets suggests that informational frictions, behavioral biases, or limits to arbitrage may prevent instantaneous price adjustment.

From a machine learning perspective, this debate is particularly important. If markets are perfectly efficient, predictive models using publicly available earnings information should not systematically outperform chance. 
Conversely, if markets exhibit delayed information incorporation, data-driven models may capture economically meaningful patterns.

Financial reports are complex, high-dimensional documents containing both structured numerical information (earnings per share, revenue growth, margins) and unstructured textual content (management discussion, risk disclosures, forward-looking statements). 
Processing this information requires cognitive effort and domain expertise. 
Behavioral finance research suggests that investors may exhibit limited attention and underreaction when confronted with complex disclosures.

For instance, Hirshleifer and Teoh\cite{limited-attention-information-disclosure} argue that investor inattention contributes to post-announcement drift. 
Similarly, Barberis, Shleifer, and Vishny\cite{model-of-investor-sentiment} develop a model of investor underreaction and overreaction driven by behavioral biases. 
These theoretical and empirical findings imply that information processing is not frictionless, creating the possibility of short-term predictability around disclosure events.

Moreover, market participants are heterogeneous. 
Institutional investors, retail traders, and algorithmic trading systems may interpret and act upon earnings information at different speeds. 
As a result, prices may adjust through a dynamic process rather than an instantaneous jump, reinforcing the empirical relevance of studying short-term post-announcement effects.

Traditional quantitative finance models frequently rely on historical price, return, and volume data. While these variables capture market dynamics and technical patterns, they do not directly encode new fundamental information. 
Earnings announcements represent discrete information shocks that can alter expectations about future cash flows and risk premia.

Empirical asset pricing research demonstrates that earnings surprises, defined as the deviation between reported earnings and market expectations, are strongly associated with abnormal returns (Bernard \& Thomas\cite{pead}). 
This suggests that numerical market data alone may be insufficient to explain short-term price adjustments following disclosures. Instead, models that incorporate financial report information may capture incremental predictive content beyond past prices.

Recent advances in machine learning and natural language processing have further emphasized the informational richness of financial text. 
For example, Loughran and Tim McDonald\cite{textual-analysis-dictionaries-and-10ks} develop domain-specific financial sentiment dictionaries, demonstrating that textual tone in corporate filings contains economically meaningful information. 
Similarly, Gu, Kelly, and Xiu\cite{empirical-asset-pricing-via-machine-learning} show that machine learning methods can extract predictive signals from large financial datasets, reshaping modern empirical asset pricing.

These developments underscore the importance of integrating structured financial data and unstructured disclosure information within predictive frameworks.

The intersection of earnings announcements and machine learning is particularly compelling. Machine learning models are well-suited to capturing nonlinear relationships, high-dimensional interactions, and complex patterns in large datasets. In financial contexts, they provide flexible tools for testing whether publicly available disclosures contain incremental predictive content beyond traditional risk factors.

However, applying machine learning to earnings-based prediction also serves a deeper purpose: it provides a data-driven test of informational efficiency. If short-term predictability persists after controlling for conventional factors, this may indicate limits to market efficiency or incomplete information diffusion. Conversely, weak predictive performance would reinforce the EMH benchmark.

Thus, studying the short-term market impact of financial reports is not merely a forecasting exercise. It directly contributes to longstanding debates in asset pricing regarding market efficiency, behavioral frictions, and the role of public disclosures in price formation.

\section{Related Work}

The intersection of machine learning and financial text analysis has evolved substantially over the past two decades. Early approaches relied on classical bag-of-words representations and linear models, while more recent work leverages deep neural networks and transformer-based architectures. In parallel, advances in representation learning have enabled increasingly sophisticated modeling of long-form financial documents such as SEC filings. This chapter reviews the most relevant strands of machine learning research related to stock price prediction using textual and market data, with particular emphasis on deep learning and domain-adapted language models.

One of the foundational works in text-based stock prediction by Tetlock\cite{tetlock2007giving} examined whether the tone of media content predicts market movements. Although not a deep learning paper, this work demonstrated that linguistic sentiment extracted from textual data contains economically meaningful signals. Textual tone was measured using dictionary-based methods, and regressions were used to link media pessimism to market returns. The key contribution was to empirically establish that textual sentiment correlates with short-term market behavior.

Subsequent research began incorporating supervised machine learning techniques. Schumaker et al.\cite{schumaker2009textual} applied support vector machines and naive Bayes classifiers to financial news articles, demonstrating that textual features could improve short-term stock direction prediction. Their objective was to classify whether a stock would increase or decrease following news releases. Using engineered linguistic features and classification algorithms, they reported predictive accuracy exceeding naive baselines, suggesting that machine learning could capture nonlinear relationships between text and returns.

These early works relied heavily on manually curated lexicons or shallow feature representations such as TF-IDF vectors. While effective to a certain degree, such representations fail to capture contextual semantics and long-range dependencies present in financial reports.


The introduction of deep learning marked a turning point in financial text modeling. Ding et al.\cite{ding2015deep} proposed a deep neural network framework that combined event embeddings derived from news headlines with temporal price modeling. Their model used a convolutional neural network (CNN) to encode event representations and incorporated historical price information. The objective was to predict stock price movements following news events. They demonstrated that deep representations significantly outperformed traditional bag-of-words approaches.

Similarly, Hu et al.\cite{hu2018listening} introduced a hierarchical attention network for modeling financial news and predicting stock volatility and direction. By applying attention mechanisms, the model could weight more informative sentences more heavily, thereby improving interpretability and performance. This work highlighted the importance of document-level structure in financial text modeling.

Parallel to developments in text modeling, recurrent neural networks (RNNs), particularly long short-term memory (LSTM) networks, were widely applied to time-series financial forecasting. Fischer et al.\cite{fischer2018deep} demonstrated that LSTMs outperform traditional autoregressive models for stock return prediction using price-based features. Although their work focused primarily on numerical time-series data rather than textual inputs, it established deep learning as a competitive alternative to classical econometric methods.

These contributions collectively show that neural architectures can effectively capture nonlinear patterns in financial data. However, many of these models were applied to relatively short textual inputs, such as headlines or summaries, rather than full-length regulatory filings.

The development of transformer architectures\cite{vaswani2017attention} and the introduction of BERT\cite{devlin2019bertpretrainingdeepbidirectional} revolutionized natural language processing. BERT's bidirectional self-attention mechanism enables contextual word representations, allowing models to capture nuanced semantics and long-range dependencies. This was particularly relevant for financial text, which is formal, domain-specific, and context-dependent.

Recognizing that general-purpose language models may not optimally capture financial terminology, researchers began developing domain-adapted transformer models. One of the most influential contributions is FinBERT\cite{yang2020finbert}, which fine-tunes BERT on financial corpora, including SEC filings and financial news. The objective of FinBERT was to improve sentiment classification accuracy in financial contexts. Empirical results showed that domain adaptation significantly improves performance compared to general BERT models when applied to financial sentiment tasks.

The existing literature demonstrates several key trends: (i) the transition from lexicon-based sentiment analysis to deep contextual embeddings, (ii) the growing importance of domain-specific language models, and (iii) the integration of textual and structured financial data within multimodal frameworks.

However, several limitations remain in prior work. First, many studies focus exclusively on news articles rather than regulatory filings, which are more comprehensive and standardized disclosures. Second, when SEC filings are used, models often truncate documents or rely on limited sections without systematic aggregation strategies. Third, relatively few works explicitly examine sector and industry-adjusted returns as alternative prediction targets, thereby limiting insight into whether models capture idiosyncratic or systematic components of stock movements.

The present study differs from prior research in several important ways. First, it focuses specifically on 10-K and 10-Q filings from U.S. markets over a recent and standardized period (2011--2024), ensuring consistency in reporting format. Second, rather than truncating long documents, it employs a controlled chunk-based sampling and aggregation strategy to approximate document-level sentiment using FinBERT. This approach balances computational feasibility with broad coverage of narrative disclosures. Third, it systematically evaluates multiple prediction horizons (1--7 days) and considers both regression and multi-class classification formulations. Finally, by constructing sector and industry-adjusted return labels using a leave-one-out methodology, this study explicitly distinguishes between firm-specific and systematic market effects.

Taken together, this work contributes to the literature by providing a comprehensive and methodologically rigorous framework for studying the short-term market impact of financial disclosures using modern transformer-based language models combined with structured price-based features.